\chapter*{Introduction}
\label{chap:introduction}
\emergencystretch=3em\relax

The Projet de Fin d'\'Etudes (PFE) closes the engineering curriculum at TELECOM
Nancy, and it is the first time the things learned separately over three years
have to hold together on a single real problem, in a real organisation, with a
deadline that belongs to someone else. Two threads of that curriculum met in
this internship. The first is software that has to be correct, where the
interesting part of the work is establishing that a program does what it claims
rather than making it run at all. The second is security understood as a
discipline of anticipating failure rather than as a list of measures to apply.
A third element was new to me: large language models, which I had used as a
tool without ever having been responsible for making one safe to put in front
of someone else. This subject offered all three at once, which is why I chose
it.

The internship took place at Wavestone, an independent management and
technology consulting group, within its Cyber practice in Luxembourg, a team of
around twenty consultants whose work includes crisis management exercises,
security audits and consulting missions for clients drawn largely from the
financial sector, the insurance sector and the European institutions. The
practice had a recurring difficulty, described at greater length in
Chapter~\ref{chap:framing}: before designing a crisis exercise or a mission for
a client, a consultant needs to know where that organisation actually stands,
and there was no structured way to establish it. The conversation opened either
on a generic checklist or on the consultant's own experience, and neither
produces something a client can audit afterwards or compare across engagements.
The mission entrusted to me, titled ``Cyber Crisis Management Maturity
Self-Assessment Platform'', set out to give that conversation a defensible
starting point.

Behind an internal need of this kind sits a broader one. In February 2024,
ENISA published \emph{Best Practices for Cyber Crisis Management}, which
organises the discipline into sixteen best practices distributed across the
four phases of a crisis, prevention, preparedness, response and recovery. It is
the European reference text on the subject, and the NIS2 directive expects
entities in essential and important sectors to be able to demonstrate a
crisis-management capability rather than simply assert one. Yet the ENISA study
is descriptive: it specifies what good crisis management looks like without
offering a scale, an evidence standard, or any way to turn sixteen qualitative
judgements into a single figure. Existing self-assessment instruments do not
close that distance either, since they were designed independently of ENISA's
structure and of NIS2's obligations. An organisation that must show how mature
it is has no purpose-built instrument with which to find out.

This is the problem the thesis addresses, and it is an engineering problem more
than a documentary one. Turning a descriptive framework into a working
instrument means answering a specific question: how can a declarative,
inherently subjective self-assessment be turned into a score that a consultant
can defend figure by figure in front of a client's security and governance
leadership, months after the fact, without the computation depending on
anything that cannot be reproduced, and without the assessment data, which is
in effect a map of an organisation's cyber weaknesses, ever leaving the
organisation's own infrastructure? The answer built during this internship
rests on three decisions. The framework itself lives in a single specification
file that no code is allowed to duplicate. A deterministic engine, written in
pure Python and holding no artificial intelligence at any stage, computes the
score from that file and the respondent's answers alone. A local language
model, running on the user's own hardware, is confined strictly downstream of
that computed result, where it writes narrative recommendations and can never
reach back into the number. Around this core sits a full web platform with real
authentication and per-owner access control. Whether that answer holds, and
what it leaves unresolved, is what the rest of this report examines.

The document is organised in nine chapters. Chapter~\ref{chap:framing}
introduces Wavestone, the Luxembourg Cyber practice, the mission itself and the
working constraints that shaped several later technical choices, before
connecting the company's initial need to the problem developed next.

Chapter~\ref{chap:context} sets out the economic, regulatory and methodological
context: cyber crises as a distinct capability from cybersecurity in general,
what NIS2 obliges organisations to demonstrate, how the ENISA framework answers
that expectation, and the precise gap that remains between the two.

Chapter~\ref{chap:state-of-the-art} reviews what already exists, general and
sector-specific maturity models, academic work on crisis-management maturity,
and the literature on using language models in sensitive domains, and states
where this thesis positions itself against all three.

Chapter~\ref{chap:method} states the approach that governed every later choice:
how the project was run sprint by sprint, the central idea of separating
deterministic scoring from AI augmentation, the five principles fixed at the
outset and never revisited, and the main alternatives that were considered and
rejected.

Chapter~\ref{chap:scoring-engine} is the core of the report. It presents the
scoring specification, the fixed-order pipeline that turns answers into a
score, the architecture of the engine itself, the testing strategy and the
independent audit it went through, and a critical reflection on the choices
within it that remain matters of judgement.

Chapter~\ref{chap:platform} describes the platform built around that engine,
the backend, the frontend, the authentication and access-control model, the
integration of the local language model, and the operational and networking
choices that kept the whole stack working across a virtual machine boundary.

Chapter~\ref{chap:validation} asks whether the result does what it claims. It
presents the three layers of testing, follows one assessment profile from raw
score to final score, evaluates the AI layer against its own guardrails, and
states the limitations that bound how much weight a result from this tool
should carry.

Chapter~\ref{chap:impact} steps back to ask what the tool is worth, to the
consulting practice, to the client organisations it targets, and to the wider
question of using language models responsibly in a regulated context, before
setting out the directions left open.

Chapter~\ref{chap:pm} turns from what was built to how it was built: the
working method in retrospect, the difficulties encountered and what each cost,
what I would do differently, and the skills the project drew on and left me
with.

The conclusion returns to the question posed above and states what was
established, what was not, and what the whole exercise taught me.